\documentclass[10pt,a4paper]{amsart}
\usepackage{amsmath,amssymb,amsthm,graphicx}
\usepackage{verbatim}
\newenvironment{reference}{\comment}{\endcomment}
%\newenvironment{reference}{}{}
\newenvironment{slogan}{\comment}{\endcomment}
\newenvironment{history}{\comment}{\endcomment}

% For commutative diagrams we use Xy-pic
\usepackage[all]{xy}

% We use 2cell for 2-commutative diagrams.
\xyoption{2cell}
% Theorem environments.
%
\theoremstyle{plain}
\newtheorem{theorem}[subsection]{Theorem}
\newtheorem{proposition}[subsection]{Proposition}
\newtheorem{lemma}[subsection]{Lemma}

\theoremstyle{definition}
\newtheorem{definition}[subsection]{Definition}
\newtheorem{example}[subsection]{Example}
\newtheorem{exercise}[subsection]{Exercise}
\newtheorem{situation}[subsection]{Situation}

\theoremstyle{remark}
\newtheorem{remark}[subsection]{Remark}
\newtheorem{remarks}[subsection]{Remarks}

\numberwithin{equation}{subsection}

% Macros
%
\def\lim{\mathop{\mathrm{lim}}\nolimits}
\def\colim{\mathop{\mathrm{colim}}\nolimits}
\def\Spec{\mathop{\mathrm{Spec}}}
\def\Hom{\mathop{\mathrm{Hom}}\nolimits}
\def\Ext{\mathop{\mathrm{Ext}}\nolimits}
\def\SheafHom{\mathop{\mathcal{H}\!\mathit{om}}\nolimits}
\def\SheafExt{\mathop{\mathcal{E}\!\mathit{xt}}\nolimits}
\def\Sch{\mathit{Sch}}
\def\Mor{\mathop{\mathrm{Mor}}\nolimits}
\def\Ob{\mathop{\mathrm{Ob}}\nolimits}
\def\Sh{\mathop{\mathit{Sh}}\nolimits}
\def\NL{\mathop{N\!L}\nolimits}
\def\CH{\mathop{\mathrm{CH}}\nolimits}
\def\proetale{{pro\text{-}\acute{e}tale}}
\def\etale{{\acute{e}tale}}
\def\QCoh{\mathit{QCoh}}
\def\Ker{\mathop{\mathrm{Ker}}}
\def\Im{\mathop{\mathrm{Im}}}
\def\Coker{\mathop{\mathrm{Coker}}}
\def\Coim{\mathop{\mathrm{Coim}}}

% Boxtimes
%
\DeclareMathSymbol{\boxtimes}{\mathbin}{AMSa}{"02}

%
% Macros for moduli stacks/spaces
%
\def\QCohstack{\mathcal{QC}\!\mathit{oh}}
\def\Cohstack{\mathcal{C}\!\mathit{oh}}
\def\Spacesstack{\mathcal{S}\!\mathit{paces}}
\def\Quotfunctor{\mathrm{Quot}}
\def\Hilbfunctor{\mathrm{Hilb}}
\def\Curvesstack{\mathcal{C}\!\mathit{urves}}
\def\Polarizedstack{\mathcal{P}\!\mathit{olarized}}
\def\Complexesstack{\mathcal{C}\!\mathit{omplexes}}
% \Pic is the operator that assigns to X its picard group, usage \Pic(X)
% \Picardstack_{X/B} denotes the Picard stack of X over B
% \Picardfunctor_{X/B} denotes the Picard functor of X over B
\def\Pic{\mathop{\mathrm{Pic}}\nolimits}
\def\Picardstack{\mathcal{P}\!\mathit{ic}}
\def\Picardfunctor{\mathrm{Pic}}
\def\Deformationcategory{\mathcal{D}\!\mathit{ef}}
\usepackage{cleveref}
\crefname{originallemma}{original lemma}{original lemmas}
\Crefname{originallemma}{Original lemma}{Original lemmas}
\makeatletter
\newlabel{lemma-sections-sheafification}{{00WK}{source}{}{sourcepointer-1}{}}
\newlabel{lemma-sections-sheafification@cref}{{[originallemma][0][]00WK}{[1][0][]source}{}{}{}}
\makeatother
\title{A map of presheaves which is bijective on a covering family becomes an isomorphism after sheafification}
\author{The Stacks Project Authors}
\date{stacks-project:sites.tex:local-snapshot-20261012; GNU FDL 1.2 or later}
\begin{document}
\maketitle
\noindent\textit{Editorial note.} All delivered mathematical text is the author's own, copied from the pinned Stacks Project chapter file sites-cohomology.tex; no statement or proof step was rewritten, reordered or removed.\par\medskip
\noindent\textit{Editorial note.} The proof is the author's complete argument; every step the author writes out is inside the counted span and no closing sentence is a placeholder.\par\medskip
\noindent\textit{Editorial note.} Every labelled reference occurring inside the counted spans is carried in assembly.external\_references as a non-load-bearing navigation pointer (same-chapter ones at their definition site, cross-chapter ones at the invocation inside the counted span), so no load-bearing mathematics is replaced by a stub. No bibliography citation occurs inside the counted spans.\par\medskip
\noindent\textit{Editorial note.} The checked chapter list for same-name collisions is the set of stacks chapters staged in this repository: chow, cohomology, divisors, etale-cohomology, flat, more-algebra, more-morphisms, morphisms, spaces-more-morphisms, varieties together with this card's chapter sites-cohomology.\par\medskip
\noindent\textit{Editorial note.} The preamble is the chapter's own retained minimal preamble (sources/stacks-sites-cohomology/retained-preamble.tex), a verbatim excerpt of the native preamble: the Xy-pic setup and the native theorem declarations with the mathematical macros.\par\medskip
\noindent\textit{Allowed standard prerequisites (editorial).}\begin{itemize}
\item Sections of a sheafification. Same chapter, Lemma lemma-sections-sheafification. Describes the sections of a sheafification by coverings and local sections, which is the tool used in both directions.
\end{itemize}\medskip
\ifdefined\hypertarget\hypertarget{sourcepointer-1}{}\fi
\noindent\textit{Source reference.} Original lemma 00WK: sources/stacks-sites/pinned-original-source.tex, source lines 2410--2410. Navigation pointer invoked at proof line 2410; the counted span contains the pointer only and the cited result is not reproduced inside it, so the pointer is registered as a non-load-bearing navigation stub in the established form for Stacks material. No bibliography citation is involved.\par\medskip
\begin{lemma}
\label{lemma-isomorphism-sheafifications}
Let $\mathcal{C}$ be a site. Let $\mathcal{F} \to \mathcal{G}$ be a map
of presheaves of sets on $\mathcal{C}$. Denote $\mathcal{B}$ the set of
$U \in \Ob(\mathcal{C})$ such that $\mathcal{F}(U) \to \mathcal{G}(U)$
is bijective. If every object of $\mathcal{C}$ has a covering by elements
of $\mathcal{B}$, then $\mathcal{F}^\# \to \mathcal{G}^\#$ is an isomorphism.
\end{lemma}

\begin{proof}
Let $U \in \Ob(\mathcal{C})$.
Let us prove that $\mathcal{F}^\#(U) \to \mathcal{G}^\#(U)$ is
surjective. To do this we will use Lemma \ref{lemma-sections-sheafification}
without further mention. For any $s \in \mathcal{G}^\#(U)$ there exists
a covering $\{U_i \to U\}$ and sections $s_i \in \mathcal{G}(U_i)$ such that
\begin{enumerate}
\item $s|_{U_i}$ is the image of $s_i$ via $\mathcal{G} \to \mathcal{G}^\#$,
and
\item for every $i, j$ there exists a covering
$\{U_{ijk} \to U_i \times_U U_j\}$ of $\mathcal{D}$ such that
the pullbacks of $s_i$ and $s_j$ to each $U_{ijk}$ agree.
\end{enumerate}
By assumption, for each $i$ we may choose a covering $\{U_{i, a} \to U_i\}$
with $U_{i, a} \in \mathcal{B}$. Then $\{U_{i, a} \to U\}$ is a covering.
Denoting $s_{i, a}$ the image of $s_i$ in $\mathcal{G}(U_{i, a})$
we see that the pullbacks of $s_{i, a}$ and $s_{j, b}$ to the members
of the covering
$$
\{(U_{i, a} \times_U U_{j, b}) \times_{U_i \times_U U_j} U_{ijk}
\to U_{i, a} \times_U U_{j, b}\}
$$
agree. Hence we may assume that $U_i \in \mathcal{B}$. Repeating
the argument, we may also assume $U_{ijk} \in \mathcal{B}$ for
all $i, j, k$ (details omitted).
Then since $\mathcal{F}(U_i) \to \mathcal{G}(U_i)$
is bijective by our definition of $\mathcal{B}$ in the statement
of the lemma, we get unique $t_i \in \mathcal{F}(U_i)$
mapping to $s_i$. The pullbacks of $t_i$ and $t_j$
to each $U_{ijk}$ agree in $\mathcal{F}(U_{ijk})$ because
$U_{ijk} \in \mathcal{B}$ and because we have the agreement
for $t_i$ and $t_j$.
Then the lemma tells us there exists a unique
$t \in \mathcal{F}^\#(U)$ such that $t|_{U_i}$ is the image of $t_i$.
It follows that $t$ maps to $s$.
We omit the proof that $\mathcal{F}^\#(U) \to \mathcal{G}^\#(U)$
is injective.
\end{proof}

\end{document}
